% from ega3-1-fr.tex
\oldpage[0\textsubscript{III}]{69}
En effet, si $\mathcal N_{hk}$ est le noyau de $u_{hk}$ pour $h\leq k$, la
suite
\[
0\longrightarrow\mathcal N_{hk}\longrightarrow\mathcal F_k
\longrightarrow\mathcal F_h\longrightarrow0
\]
est exacte d'après (iii), et notre assertion résulte de ce que le foncteur
$\mathcal A\mapsto\mathcal C^i(X,\mathcal A)$ est exact (G, II, 4.3).
Soit
$\mathcal G^i=\varprojlim_k\mathcal C^i(X,\mathcal F_k)$~; on a donc
$H^j(X,\mathcal G^i)=0$ pour $j>0$ et $i\geq0$ en vertu de \emph{a)}. Nous
allons montrer que $\mathcal G^\bullet=(\mathcal G^i)_{i\geq0}$ est une
résolution du faisceau $\mathcal F$~; comme
$H^j(X,\mathcal G^i)=0$ pour $j>0$, la cohomologie $H^\bullet(X,\mathcal F)$
sera égale à $H^\bullet(\Gamma(X,\mathcal G^\bullet))$ (G, II, 4.7.1).

Il est clair que, par passage à la limite projective, on déduit des suites
exactes
\[
0\longrightarrow\mathcal F_k
\longrightarrow\mathcal C^0(X,\mathcal F_k)
\longrightarrow\mathcal C^1(X,\mathcal F_k)\longrightarrow\cdots
\]
un complexe de faisceaux de groupes abéliens
\[
0\longrightarrow\mathcal F\longrightarrow\mathcal G^0
\longrightarrow\mathcal G^1\longrightarrow\cdots.
\]
Pour prouver notre assertion, il faut établir que
$\mathcal H^i(\mathcal G^\bullet)=0$ pour $i>0$. Ce faisceau est engendré
par le préfaisceau
$U\mapsto H^i(\Gamma(U,\mathcal G^\bullet))$ (G, II, 4.1)~; or, le complexe
$\Gamma(U,\mathcal G^\bullet)$ est la limite projective du système
projectif de complexes de groupes abéliens
\[
(\Gamma(U,\mathcal C^\bullet(X,\mathcal F_k)))_{k\in\mathbb N}
\]
(0\textsubscript{I}, 3.2.6). On a vu dans \emph{a)} que pour chaque
$i\geq0$, les applications
$\Gamma(U,\mathcal C^i(X,\mathcal F_k))\to
\Gamma(U,\mathcal C^i(X,\mathcal F_h))$ $(h\leq k)$ sont surjectives~;
d'autre part, on a
$H^i(U,\mathcal F_k)=
H^i(\Gamma(U,\mathcal C^\bullet(X,\mathcal F_k)))$, la résolution canonique
$\mathcal C^\bullet(U,\mathcal F_k|U)$ étant induite sur $U$ par
$\mathcal C^\bullet(X,\mathcal F_k)$~; en vertu de l'hypothèse (i), pour
tout $U\in\mathfrak B$, on peut appliquer
(\hyperref[0.13.2.3-fr]{13.2.3}) au système projectif de complexes
$(\Gamma(U,\mathcal C^\bullet(X,\mathcal F_k)))_{k\in\mathbb N}$, et on a
donc
\[
H^i(\Gamma(U,\mathcal G^\bullet))
=\varprojlim_k H^i(U,\mathcal F_k)
\quad\hbox{pour tout }i\geq0.
\]
L'hypothèse (ii) prouve bien alors, par définition, que les faisceaux
$\mathcal H^i(\mathcal G^\bullet)$ sont nuls pour $i>0$.

On a alors pour tout $i\geq0$,
\[
H^i(X,\mathcal F)=H^i(\Gamma(X,\mathcal G^\bullet))
\quad\hbox{et}\quad
\Gamma(X,\mathcal G^\bullet)
=\varprojlim_k\Gamma(X,\mathcal C^\bullet(X,\mathcal F_k)).
\]
On vient de remarquer que les applications
$\Gamma(X,\mathcal C^i(X,\mathcal F_k))\to
\Gamma(X,\mathcal C^i(X,\mathcal F_h))$ $(h\leq k)$ sont toutes
surjectives~; la conclusion résulte donc encore de
(\hyperref[0.13.2.3-fr]{13.2.3}).
\end{proof}

\begin{env}[13.3.2]
\label{0.13.3.2-fr}
\emph{Remarques} (13.3.2).
\begin{enumerate}
\item[\emph{(i)}]
L'énoncé (\hyperref[0.13.3.1-fr]{13.3.1}) n'a d'intérêt que pour $i>0$,
puisque pour $i=0$, $h_i$ est toujours un isomorphisme sans hypothèse
(0\textsubscript{I}, 3.2.6).
\item[\emph{(ii)}]
Les conditions (i) et (ii) de (\hyperref[0.13.3.1-fr]{13.3.1}) seront en
particulier vérifiées si $H^i(U,\mathcal F_k)=0$ pour tout $k$, tout $i>0$
et tout $U\in\mathfrak B$, et si pour $U\in\mathfrak B$, les applications
$\Gamma(U,\mathcal F_k)\to\Gamma(U,\mathcal F_h)$ sont surjectives. Ce sera
le cas le plus fréquent d'application de
(\hyperref[0.13.3.1-fr]{13.3.1}).
\end{enumerate}
\end{env}

\subsection{Condition de Mittag-Leffler et objets gradués associés aux systèmes projectifs}
\label{subsection:0.13.4-fr}
\label{0.13.4-fr}

\begin{env}[13.4.1]
\label{0.13.4.1-fr}
Soit $\mathbf A=(A_k,u_{kh})_{k\in\mathbb Z}$ un système projectif dans une
catégorie abélienne $C$~; nous dirons qu'il est \emph{limité inférieurement}
s'il existe $k_0$ tel que $A_k=0$ pour $k<k_0$. Nous définirons sur chaque
$A_k$ une \emph{filtration} $(\mathrm F^p(A_k))_{p\in\mathbb Z}$ par les
formules
\begin{equation}
\label{0.13.4.1.1-fr}
\tag{13.4.1.1}
\mathrm F^p(A_k)=
\begin{cases}
\operatorname{Ker}(A_k\longrightarrow A_{p-1})&\text{pour }p\leq k+1,\\
0&\text{pour }p\geq k+1.
\end{cases}
\end{equation}
